\slajdid{04-s033}
\begin{frame}[label=04-s033]{Komplement maksimalne vrednosti i komplement osnove}
\gusttekst
Ova dva pojma se uvode u nedostatku simbola, znaka, minus kako bi se prikazali i negativni brojevi, a koji će omogućiti da se na „lak“ način izvode aritmetičke operacije i sa pozitivnim i sa negativnim brojevima.

U opštem slučaju za brojni sistem sa osnovom $r$ i prikazom sa dužinom reči $n$ (broj cifara)\\
\alert{komplement maksimalne vrednosti}, odnosno prikaz negativnog broja je
\[-V\equiv(r^n-1)-V\]
U binarnom brojnom sistemu – komplement jedinice ili prvi komplement

Zanimljiv pošto se dobija direktnim invertovanjem vrednosti bita
\[111\ldots11-b_{n-1}b_{n-2}\,b_{n-3}\ldots b_1\,b_0=\overline{b_{n-1}}\,\overline{b_{n-2}}\,\overline{b_{n-3}}\ldots\overline{b_1}\,\overline{b_0}\]
PRIMER: u decimalnom brojnom sistemu sa tri cifre broj različitih vrednosti je 1000. Ako vrednosti od 0 do 499 smatramo pozitivnim onda će negativna vrednost -123 biti prikazana u komplementu 9tke kao
\[-123\equiv999-123=876\]
Odnosno vodeće cifre 0,1,2,3,4 će biti za pozitivne brojeve a 5,6,7,8,9 za negativne brojeve
\end{frame}
